\documentclass[a4paper,10pt]{report}\input{../0.Préambule/Préambule_cours_prof}\begin{document}

\newpage
\section*{Puissance de dix}
\addcontentsline{toc}{section}{Puissance de dix}

\subsection*{Découvrir les puissances de dix}
\addcontentsline{toc}{subsection}{Découvrir les puissances de dix}

\begin{exer1}
Écris sous forme décimale les puissances de $10$ suivantes.

\begin{enumerate}[font=\color{exer1}]
\begin{tabularx}{\linewidth}{*{8}{>{\arraybackslash}X}}
\item $10^2$ &
\item $10^6$ &
\item $10^{-2}$ &
\item $10^0$ &
\item $10^{-4}$ &
\item $10^{-1}$ &
\item $10^1$ &
\item $10^5$ \\
\end{tabularx}
\end{enumerate}
\end{exer1}

\noindent \begin{minipage}{8.5cm}
\begin{exer1}
Écris sous forme d'une puissance de $10$.

\begin{enumerate}[font=\color{exer1}]
\begin{tabularx}{\linewidth}{*{3}{>{\arraybackslash}X}}
\item $\np{10000}$ &
\item $\np{100000}$ &
\item $\np{0,01}$ \\
\item $\np{0,001}$ &
\item $\np{0,000001}$ &
\item $1$ \\
\end{tabularx}
\end{enumerate}
\end{exer1}

\begin{exer1}
Écris sous forme d'une puissance de $10$.

\begin{enumerate}[font=\color{exer1}]
\begin{tabularx}{\linewidth}{*{3}{>{\arraybackslash}X}}
\item $\np{1000000}$ &
\item $\np{10000000}$ &
\item $\np{0,0000001}$ \\
\item $\np{0,1}$ &
\item $\np{0,00001}$ &
\item $10$ \\
\end{tabularx}
\end{enumerate}
\end{exer1}
\end{minipage}
\hspace{10mm}
\begin{minipage}{7cm}
\begin{exer1}
Relie les expressions égales.

\renewcommand{\arraystretch}{1.5}
\begin{tabularx}{7cm}{r c X c l}
$\np{100000000}$ & $\bullet$ & & $\bullet$ & $10^{-6}$ \\
$\np{0,000001}$ & $\bullet$ & & $\bullet$ & $10^{-9}$ \\
$\np{10000}$ & $\bullet$ & & $\bullet$ & $10^{-2}$ \\
$\np{0,01}$ & $\bullet$ & & $\bullet$ & $10^{4}$ \\
$\np{0,001}$ & $\bullet$ & & $\bullet$ & $10^{8}$ \\
$\np{0,000000001}$ & $\bullet$ & & $\bullet$ & $10^{-3}$ \\
\end{tabularx}
\end{exer1}
\end{minipage}

\begin{exer1}
Exprime les grandeurs suivantes sous la forme d'une puissance de $10$.

\begin{enumerate}[font=\color{exer1}]
\begin{tabularx}{\linewidth}{*{3}{>{\arraybackslash}X}}
\item $1$cm = \hspace{1.5cm} m &
\item $1$dam = \hspace{1.5cm} m &
\item $1$km = \hspace{1.5cm} m \\
\item $1$mm = \hspace{1.5cm} m &
\item $1$hm = \hspace{1.5cm} m &
\item $1$dm = \hspace{1.5cm} m \\
\end{tabularx}
\end{enumerate}
\end{exer1}

\begin{exer1}
Exprime les grandeurs suivantes sous la forme d'une puissance de $10$.

\begin{enumerate}[font=\color{exer1}]
\begin{tabularx}{\linewidth}{*{4}{>{\arraybackslash}X}}
\item $\np{10}$cm = \hspace{1cm} mm &
\item $\np{100}$dam = \hspace{1cm} m &
\item $\np{10}$km = \hspace{1cm} dam &
\item $\np{100}$mm = \hspace{1cm} m \\
\item $\np{10}$dm = \hspace{1cm} km &
\item $\np{100000}$dm = \hspace{1cm} m &
\item $\np{1000}$m = \hspace{1cm} km &
\item $\np{10000}$cm = \hspace{1cm} m \\
\end{tabularx}
\end{enumerate}
\end{exer1}

\begin{exer1}
Exprime les grandeurs suivantes sous la forme d'une puissance de $10$.

\begin{enumerate}[font=\color{exer1}]
\begin{tabularx}{\linewidth}{*{3}{>{\arraybackslash}X}}
\item $\np{10}$g = \hspace{1.5cm} kg &
\item $\np{10}$Mo = \hspace{1.5cm} o &
\item $\np{1}$mg = \hspace{1.5cm} hg \\
\item $\np{1000}$hL = \hspace{1.5cm} L &
\item $\np{1}$ng = \hspace{1.5cm} g &
\item $\np{1}$ $\mu$m = \hspace{1.5cm} m \\
\end{tabularx}
\end{enumerate}
\end{exer1}

\subsection*{Calculer avec les puissances de dix}
\addcontentsline{toc}{subsection}{Calculer avec les puissances de dix}

\begin{exer1}
Écris sous forme décimale puis d'une puissance de $10$.

\begin{enumerate}[font=\color{exer1}]
\begin{tabularx}{\linewidth}{*{6}{>{\arraybackslash}X}}
\item $10^2 \times 10^4$ =&
\item $10^5 \times 10^2$ =&
\item $10^{-2} \times 10^3$ =&
\item $10^{-1} \times 10^4$ =&
\item $10^1 \times 10^{-2}$ =&
\item $10^{-5} \times 10^{-3}$ =\\
\end{tabularx}
\end{enumerate}
\end{exer1}

\noindent \begin{minipage}{7.8cm}
\begin{exer1}
Complète par une puissance de $10$.

\renewcommand{\arraystretch}{1.5}
\begin{tabularx}{7cm}{|*{5}{>{\centering \arraybackslash}X|}}
\hline
\rowcolor{exer1!30} $\times$ & $10^{9}$ & $10^{-7}$ & $10^{56}$ & $10^{-18}$ \\\hline
\cellcolor{exer1!30}{$10^{12}$} & $10^{21}$ & & & \\\hline
\cellcolor{exer1!30}{$10^{9}$} & & & & \\\hline
\cellcolor{exer1!30}{$10^{-25}$} & & & & \\\hline
\cellcolor{exer1!30}{$10^{8}$} & & & & \\\hline
\end{tabularx}
\end{exer1}
\end{minipage}
\hspace{5mm}
\begin{minipage}{8.5cm}
\begin{exer1}
Écris sous forme décimale puis d'une puissance de $10$.

\vspace{-2mm}
\begin{enumerate}[font=\color{exer1}]
\begin{tabularx}{\linewidth}{*{4}{>{\arraybackslash}X}}
\item ${\big( 10^{4} \big)}^{2}$ =&
\item ${\big( 10^{2} \big)}^{3}$ =&
\item ${\big( 10^{1} \big)}^{5}$ =&
\item ${\big( 10^{-2} \big)}^{4}$ =\\
\end{tabularx}
\end{enumerate}
\end{exer1}

\begin{exer1}
Écris sous forme décimale puis d'une puissance de $10$.

\vspace{-2mm}
\begin{enumerate}[font=\color{exer1}]
\begin{tabularx}{\linewidth}{*{4}{>{\arraybackslash}X}}
\item $\dfrac{10^{7}}{10^{4}}$ =&
\item $\dfrac{10^{4}}{10^{3}}$ =&
\item $\dfrac{10^{12}}{10^{9}}$ =&
\item $\dfrac{10^{2}}{10^{3}}$ =\\
\end{tabularx}
\end{enumerate}
\end{exer1}
\end{minipage}

\begin{exer2}
Calcule les expressions suivantes.

\vspace{-2mm}
\begin{enumerate}[font=\color{exer2}]
\begin{tabularx}{\linewidth}{p{4cm}*{3}{>{\arraybackslash}X}}
\item \phantom{$\dfrac{1}{2}$} $10^{5} + 10^{2} + 10^{5} \times 10^{3}$ =&
\item \phantom{$\dfrac{1}{2}$} $10^{4} - 10^{3} + 10^{-1}$ =&
\item $\dfrac{10^{4}-10^{3}-10^{2}}{{\big( 10^{2} \big)}^{2}}$ =&
\item $\dfrac{10^{4}+10^{3}+10^{2}}{10^{2}}$ =\\
\end{tabularx}
\end{enumerate}
\end{exer2}

\subsection*{Écrire un nombre en notation scientifique}
\addcontentsline{toc}{subsection}{Écrire un nombre en notation scientifique}

\vspace{-2mm}
\begin{exer1}
Donne l'écriture décimale de chaque nombre.

\vspace{-1mm}
\begin{enumerate}[font=\color{exer1}]
\begin{tabularx}{\linewidth}{*{6}{>{\arraybackslash}X}}
\item $1,35 \times 10^5$ &
\item $\np{0,00605} \times 10^2$ &
\item $\np{45200} \times 10^{-5}$ &
\item $\np{2} \times 10^{-4}$ &
\item $\np{0,05} \times 10^4$ &
\item $\np{13,45} \times 10^{-3}$ \\
\end{tabularx}
\end{enumerate}
\end{exer1}

\begin{exer2}
Complète.

\vspace{-1mm}
\begin{enumerate}[font=\color{exer2}]
\begin{tabularx}{\linewidth}{*{3}{>{\arraybackslash}X}}
\item $1,45 \times 10 \quad \quad = \np{14500}$ &
\item $45 \times 10 \quad \quad = \np{0,0045}$ &
\item $-6,3 \times 10 \quad \quad = \np{-6300}$ \\
\item $\quad \quad \quad \times 10^{-2} = \np{85}$ &
\item $\quad \quad \quad \times 10^6 = \np{7,1}$ &
\item $\quad \quad \quad \times 10^{-3} = \np{-0,063}$ \\
\end{tabularx}
\end{enumerate}
\end{exer2}

\begin{exer3}
Dans chaque cas, déterminer la valeur de $n$ ou de $x$ manquante vérifiant l'égalité :

\vspace{-1mm}
\begin{enumerate}[font=\color{exer3}]
\begin{tabularx}{\linewidth}{*{3}{>{\arraybackslash}X}}
\item $532\times 10^{n} = 5,32$ &
\item $67\times 10^{n} = 0,00067$ & 
\item $x\times 10^{3} = 531,8$ \\
\item $6,54\times 10^{5} = 654\times 10^{n}$ &
\item $6,12\times 10^{-13} = x\times 10^{-12}$ &
\item $0,561\times 10^{-7} = 56,1 \times 10^n$ \\
\end{tabularx}
\end{enumerate}
\end{exer3}

\begin{exer1}
Écris chaque nombre en notation scientifique.

\vspace{-1mm}
\begin{enumerate}[font=\color{exer1}]
\begin{tabularx}{\linewidth}{p{2cm}p{2.5cm}*{2}{>{\arraybackslash}X}}
\item deux-mille &
\item cinq millions &
\item quarante-sept millièmes &
\item cinquante-deux millionièmes \\
\end{tabularx}
\end{enumerate}
\end{exer1}

\begin{exer2}
Écris chaque nombre sous forme décimale puis en notation scientifique.

\vspace{-1mm}
\begin{enumerate}[font=\color{exer2}]
\begin{tabularx}{\linewidth}{*{2}{>{\arraybackslash}X}}
\item quatre-mille-cinq-cent-trois &
\item huit-cent-mille-quatre-vingt-douze \\
\item deux-millions et douze-millièmes &
\item trente-neuf-millièmes et quarante-dixièmes \\
\item soixante-dix-huit-millionièmes &
\item cent-mille-million de milliards \\
\end{tabularx}
\end{enumerate}
\end{exer2}

\begin{exer1}
Entoure les nombres qui sont écrit en notation scientifique.

\vspace{-1mm}
\begin{enumerate}[font=\color{exer1}]
\begin{tabularx}{\linewidth}{*{5}{>{\arraybackslash}X}}
\item[]$3,4\times 10^{-2}$	&
\item[]$45\times 10^{3}$ &
\item[]$0,95\times 10^{4}$ &
\item[]$4,7\times 10^{-2}$&
\item[]$10,1\times 10^{3}$ \\
\item[]$4,3\times 2^{4}$ &
\item[]$3,7\times 10^{14}$ &
\item[]$8,76\times 10^{17}$&
%\item[]$0,21\times 10^{4}$ &
\item[]$4,25\times 10^{37}$ &
\item[]$321\times 10^{3}$ \\
\end{tabularx}
\end{enumerate}
\end{exer1}

\begin{exer1}
Donner la notation scientifique des nombres suivants.

\vspace{-1mm}
\begin{enumerate}[font=\color{exer1}]
\begin{tabularx}{\linewidth}{*{4}{>{\arraybackslash}X}}
\item $\np{632 200}$ = \dots ~ \dots   ~ \dots &
\item $\np{0,000 032}$ = \dots ~ \dots   ~ \dots &
\item $\np{0,000 000 05}$ =  \dots ~ \dots   ~ \dots &
\item $\np{6 540 000}$ = \dots ~ \dots   ~ \dots \\
\end{tabularx}
\end{enumerate}
\end{exer1}

\begin{exer2}
Donner les écritures scientifiques des nombres ci-dessous :

\vspace{-1mm}
\begin{enumerate}[font=\color{exer2}]
\begin{tabularx}{\linewidth}{*{6}{>{\arraybackslash}X}}
\item $\np{4 540 000}$&
\item $\np{0,000 054}$&
\item $-354,1\times 10^{11}$&
\item $79,8\times 10^{-8}$&
\item $-0,0079\times 10^{8}$&
\item $0,0042\times 10^{-4}$\\
\end{tabularx}
\end{enumerate}
\end{exer2}

\begin{exer2}
On a prélevé$ 1$mL de sang d'un adulte. Dans cet échantillon, il y a $43 \times 10^5$ globules rouges. Le corps de cet adulte contient $5$L de sang.

\medskip \noindent Combien de globules rouges contient le corps de cette personne ? On donnera la réponse en écriture scientifique.
\end{exer2}

\begin{exer2}
La vitesse d'une sonde spatiale est d'environ $\np{20 800}$ m/s.
Donner l'écriture scientifique de cette vitesse exprimée en kilomètres par heure.
\end{exer2}

\begin{exer3}
La masse d'un atome de cuivre est de $1,05 \times 10^{-30}$g. Combien d'atomes de cuivre y a t-il dans $1,47$kg de cuivre ?
\end{exer3}

\begin{exer3}
$1$m$^3$ d'eau de mer contient $0,004$mg d'or. Sur la Terre, le volume total d'eau de mer est d'environ $1,3 \times 10^{6}$km$^3$. Calcule la masse totale d'or en tonnes que renferment les mers et les océans sur Terre. Écris le résultat en notation scientifique.
\end{exer3}

\begin{exer3}
La structure de la tour Eiffel a une masse de $\np{7300}$ tonnes. On considère que la structure est composée essentiellement de fer. Sachant qu'un atome de fer a une masse de $9,352 \times 10^{-26}$kg, combien d'atome de fer y a-t-il dans la structure ? Donne une valeur arrondie à l'unité.
\end{exer3}
\end{document}